\section{Original ample cuts in the numerical ampleness induction}

The hypothesis on lower-dimensional closed subschemes in this section is
an intermediate dimension-induction conclusion. It is used only to prove
the displayed mixed inequalities on the original ambient scheme; the
final nef-plus-ample theorem must supply this hypothesis internally.

\begin{lemma}[Strict first ample entry from lower-dimensional ampleness {\cite[\S1.2]{BJ22}}]
\label{lem:native-nef-ample-first-cut-positive}
\lean{EnvelopeProperty.SeshadriAmplenessCriterionNefAmpleCutPositivity.mixedEulerDifference_ample_cons_pos_of_closed_pullback_ample}
\leanok
\uses{thm:seshadri-ample-top-intersection-positive,lem:seshadri-mixed-euler-section-cut,lem:intersection-power-entry,thm:seshadri-mixed-euler-dimension-vanishing,thm:projective-original-ample-section-prescribed-closed-set,lem:seshadri-original-section-dimension-drop}
Let $p:X\to\operatorname{Spec}k$ be projective with $X$ integral
of dimension $d+1$. Let $A$ be ample and let $B$ be invertible.
Suppose every original closed subscheme $i:Y\hookrightarrow X$
of dimension at most $d$ has original $i^*B$ ample. Then
\[
 \Delta_A\Delta_B^d\chi(X,\mathcal O_X)>0.
\]
\end{lemma}
\begin{proof}
Choose a literal reduced integral point closure $Z\hookrightarrow X$
of dimension $d$. It is a proper closed subset of the integral
$X$. Ampleness produces a nonzero section $s$ of $A^r$, with
$r>0$, whose original zero scheme $i:D\hookrightarrow X$ contains
this same $Z$ by an actual closed immersion $j:Z\hookrightarrow D$.
The actual Cartier dimension-drop theorem gives $\dim D\leq d$.
The induction hypothesis therefore makes the original $B_D=i^*B$
ample, and closed restriction makes the original $j^*B_D$ ample.

Strict top self-intersection positivity on the integral $Z$ gives
$\Delta_{j^*B_D}^d\chi(Z,\mathcal O_Z)>0$. Monotonicity of top
ample Euler differences under a closed immersion into $D$ bounds
this value above by $\Delta_{B_D}^d\chi(D,\mathcal O_D)$.
Thus the latter value is positive, even though $D$ itself may be
reducible or nonreduced.

The actual section exact sequence gives
\[
 r\,\Delta_A\Delta_B^d\chi(X,\mathcal O_X)
 =\Delta_{B_D}^d\chi(D,i^*A^r).
\]
Top mixed coefficient independence identifies its right side with
the positive value having coefficient $\mathcal O_D$. Divide by
the positive integer $r$ to obtain the required strict inequality.
\end{proof}

\begin{lemma}[Two ample entries from lower-dimensional ampleness {\cite[\S1.2]{BJ22}}]
\label{lem:native-nef-ample-second-cut-nonnegative}
\lean{EnvelopeProperty.SeshadriAmplenessCriterionNefAmpleCutPositivity.mixedEulerDifference_two_ample_cons_nonneg_of_closed_pullback_ample}
\leanok
\uses{thm:native-nef-ample-mixed-euler-nonnegative,thm:native-numerical-ample-classes,thm:seshadri-ample-top-intersection-positive,lem:seshadri-mixed-euler-section-cut,lem:intersection-power-entry,thm:projective-original-ample-section-prescribed-closed-set,lem:seshadri-original-section-dimension-drop}
Let $X$ be integral and projective with $\dim X\leq d+2$, let
$A$ be ample and $B$ invertible. Suppose every original closed
subscheme of dimension at most $d+1$ has original pullback of $B$
ample. Then
\[
 \Delta_A^2\Delta_B^d\chi(X,\mathcal O_X)\geq0.
\]
\end{lemma}
\begin{proof}
Choose a nonzero section of an actual positive ample power $A^r$,
and form its original Cartier zero scheme $i:D\hookrightarrow X$.
It has dimension at most $d+1$. By the induction hypothesis
$B_D=i^*B$ is ample, while $A_D=i^*A$ is ample by actual closed
restriction and hence is nef for all original closed-curve tests.
The preceding mixed nef--ample nonnegativity theorem, with the
actual invertible coefficient $i^*A^r$, gives
\[
  \Delta_{A_D}\Delta_{B_D}^d\chi(D,i^*A^r)\geq0.
\]
The section exact sequence and top tensor-power identity identify
this value with
$r\Delta_A^2\Delta_B^d\chi(X,\mathcal O_X)$.
Since $r>0$, the desired nonnegativity follows. No ampleness of
$B$ on the ambient $X$ has been assumed.
\end{proof}
